% Front-end state of a multithreaded SM (Lesson 4): one program counter
% per warp (a PC table W0..W3), feeding the instruction cache and decoder,
% and a register file with a private register set for every thread,
% grouped by warp.  Redraws the outline figure m4l1_fe_extension.
% Usage: % Front-end state of a multithreaded SM (Lesson 4): one program counter
% per warp (a PC table W0..W3), feeding the instruction cache and decoder,
% and a register file with a private register set for every thread,
% grouped by warp.  Redraws the outline figure m4l1_fe_extension.
% Usage: % Front-end state of a multithreaded SM (Lesson 4): one program counter
% per warp (a PC table W0..W3), feeding the instruction cache and decoder,
% and a register file with a private register set for every thread,
% grouped by warp.  Redraws the outline figure m4l1_fe_extension.
% Usage: % Front-end state of a multithreaded SM (Lesson 4): one program counter
% per warp (a PC table W0..W3), feeding the instruction cache and decoder,
% and a register file with a private register set for every thread,
% grouped by warp.  Redraws the outline figure m4l1_fe_extension.
% Usage: \input{figures/l04-warp-state}
% Width: body (120 mm); the drawing is about 118 mm wide.
\begin{tikzpicture}[x=1mm, y=1mm, font=\scriptsize\sffamily,
    >={Stealth[length=1.6mm]},
    blk/.style={draw, fill=ln-box, minimum height=6mm, inner sep=2pt,
                align=center}]
  % PC table
  \node[anchor=south west, font=\scriptsize\sffamily\bfseries] at (0,27)
    {\iflangja{ワープごとの PC}{one PC per warp}};
  \foreach \w/\pc [count=\i from 0] in {W0/0x8000, W1/0x8004, W2/0x8000, W3/0x800C} {
    \draw[fill=ln-accent!12] (0,26-\i*6) rectangle ++(8,-5.5)
      node[midway] {\w};
    \draw[fill=white] (8,26-\i*6) rectangle ++(14,-5.5)
      node[midway, font=\scriptsize\ttfamily] {\pc};
  }
  % fetch path
  \node[blk, minimum width=14mm] (icache) at (36,14)
    {\iflangja{命令\\キャッシュ}{I-cache}};
  \node[blk, minimum width=14mm, right=3mm of icache] (dec)
    {\iflangja{デコーダ}{Decoder}};
  \draw[->] (22,14) -- (icache);
  \draw[->] (icache) -- (dec);
  \node[anchor=north, text=ln-muted, align=center] at (36,8.5)
    {\iflangja{選ばれたワープの\\PC で1命令をフェッチ}{one instruction fetched\\at the selected warp's PC}};
  % register file: 4 warps x 4 threads, registers R0..R3
  \node[anchor=south west, font=\scriptsize\sffamily\bfseries] at (66,27)
    {\iflangja{スレッドごとのレジスタ}{a register set per thread}};
  \foreach \w in {0,...,3} {
    \draw[fill=ln-accent!12] (66+\w*13,26) rectangle ++(12.6,-4)
      node[midway] {W\w};
    \foreach \t in {0,...,3} {
      \node[font=\scriptsize\sffamily] at (66+\w*13+\t*3.2+1.6,20.2) {T\t};
    }
  }
  \foreach \r in {0,...,3} {
    \node[anchor=east, inner sep=1pt] at (65.5,16.6-\r*3.4) {R\r};
    \foreach \w in {0,...,3} {
      \foreach \t in {0,...,3} {
        \draw[fill=white] (66+\w*13+\t*3.2,18.3-\r*3.4) rectangle ++(3.2,-3.4);
      }
    }
  }
  \node[anchor=north west, text=ln-muted, align=left] at (66,3)
    {\iflangja{W0 の R1 読み出し = T0--T3 の R1 を同時に読む}{reading R1 for W0 = R1 of T0--T3 at once}};
\end{tikzpicture}

% Width: body (120 mm); the drawing is about 118 mm wide.
\begin{tikzpicture}[x=1mm, y=1mm, font=\scriptsize\sffamily,
    >={Stealth[length=1.6mm]},
    blk/.style={draw, fill=ln-box, minimum height=6mm, inner sep=2pt,
                align=center}]
  % PC table
  \node[anchor=south west, font=\scriptsize\sffamily\bfseries] at (0,27)
    {\iflangja{ワープごとの PC}{one PC per warp}};
  \foreach \w/\pc [count=\i from 0] in {W0/0x8000, W1/0x8004, W2/0x8000, W3/0x800C} {
    \draw[fill=ln-accent!12] (0,26-\i*6) rectangle ++(8,-5.5)
      node[midway] {\w};
    \draw[fill=white] (8,26-\i*6) rectangle ++(14,-5.5)
      node[midway, font=\scriptsize\ttfamily] {\pc};
  }
  % fetch path
  \node[blk, minimum width=14mm] (icache) at (36,14)
    {\iflangja{命令\\キャッシュ}{I-cache}};
  \node[blk, minimum width=14mm, right=3mm of icache] (dec)
    {\iflangja{デコーダ}{Decoder}};
  \draw[->] (22,14) -- (icache);
  \draw[->] (icache) -- (dec);
  \node[anchor=north, text=ln-muted, align=center] at (36,8.5)
    {\iflangja{選ばれたワープの\\PC で1命令をフェッチ}{one instruction fetched\\at the selected warp's PC}};
  % register file: 4 warps x 4 threads, registers R0..R3
  \node[anchor=south west, font=\scriptsize\sffamily\bfseries] at (66,27)
    {\iflangja{スレッドごとのレジスタ}{a register set per thread}};
  \foreach \w in {0,...,3} {
    \draw[fill=ln-accent!12] (66+\w*13,26) rectangle ++(12.6,-4)
      node[midway] {W\w};
    \foreach \t in {0,...,3} {
      \node[font=\scriptsize\sffamily] at (66+\w*13+\t*3.2+1.6,20.2) {T\t};
    }
  }
  \foreach \r in {0,...,3} {
    \node[anchor=east, inner sep=1pt] at (65.5,16.6-\r*3.4) {R\r};
    \foreach \w in {0,...,3} {
      \foreach \t in {0,...,3} {
        \draw[fill=white] (66+\w*13+\t*3.2,18.3-\r*3.4) rectangle ++(3.2,-3.4);
      }
    }
  }
  \node[anchor=north west, text=ln-muted, align=left] at (66,3)
    {\iflangja{W0 の R1 読み出し = T0--T3 の R1 を同時に読む}{reading R1 for W0 = R1 of T0--T3 at once}};
\end{tikzpicture}

% Width: body (120 mm); the drawing is about 118 mm wide.
\begin{tikzpicture}[x=1mm, y=1mm, font=\scriptsize\sffamily,
    >={Stealth[length=1.6mm]},
    blk/.style={draw, fill=ln-box, minimum height=6mm, inner sep=2pt,
                align=center}]
  % PC table
  \node[anchor=south west, font=\scriptsize\sffamily\bfseries] at (0,27)
    {\iflangja{ワープごとの PC}{one PC per warp}};
  \foreach \w/\pc [count=\i from 0] in {W0/0x8000, W1/0x8004, W2/0x8000, W3/0x800C} {
    \draw[fill=ln-accent!12] (0,26-\i*6) rectangle ++(8,-5.5)
      node[midway] {\w};
    \draw[fill=white] (8,26-\i*6) rectangle ++(14,-5.5)
      node[midway, font=\scriptsize\ttfamily] {\pc};
  }
  % fetch path
  \node[blk, minimum width=14mm] (icache) at (36,14)
    {\iflangja{命令\\キャッシュ}{I-cache}};
  \node[blk, minimum width=14mm, right=3mm of icache] (dec)
    {\iflangja{デコーダ}{Decoder}};
  \draw[->] (22,14) -- (icache);
  \draw[->] (icache) -- (dec);
  \node[anchor=north, text=ln-muted, align=center] at (36,8.5)
    {\iflangja{選ばれたワープの\\PC で1命令をフェッチ}{one instruction fetched\\at the selected warp's PC}};
  % register file: 4 warps x 4 threads, registers R0..R3
  \node[anchor=south west, font=\scriptsize\sffamily\bfseries] at (66,27)
    {\iflangja{スレッドごとのレジスタ}{a register set per thread}};
  \foreach \w in {0,...,3} {
    \draw[fill=ln-accent!12] (66+\w*13,26) rectangle ++(12.6,-4)
      node[midway] {W\w};
    \foreach \t in {0,...,3} {
      \node[font=\scriptsize\sffamily] at (66+\w*13+\t*3.2+1.6,20.2) {T\t};
    }
  }
  \foreach \r in {0,...,3} {
    \node[anchor=east, inner sep=1pt] at (65.5,16.6-\r*3.4) {R\r};
    \foreach \w in {0,...,3} {
      \foreach \t in {0,...,3} {
        \draw[fill=white] (66+\w*13+\t*3.2,18.3-\r*3.4) rectangle ++(3.2,-3.4);
      }
    }
  }
  \node[anchor=north west, text=ln-muted, align=left] at (66,3)
    {\iflangja{W0 の R1 読み出し = T0--T3 の R1 を同時に読む}{reading R1 for W0 = R1 of T0--T3 at once}};
\end{tikzpicture}

% Width: body (120 mm); the drawing is about 118 mm wide.
\begin{tikzpicture}[x=1mm, y=1mm, font=\scriptsize\sffamily,
    >={Stealth[length=1.6mm]},
    blk/.style={draw, fill=ln-box, minimum height=6mm, inner sep=2pt,
                align=center}]
  % PC table
  \node[anchor=south west, font=\scriptsize\sffamily\bfseries] at (0,27)
    {\iflangja{ワープごとの PC}{one PC per warp}};
  \foreach \w/\pc [count=\i from 0] in {W0/0x8000, W1/0x8004, W2/0x8000, W3/0x800C} {
    \draw[fill=ln-accent!12] (0,26-\i*6) rectangle ++(8,-5.5)
      node[midway] {\w};
    \draw[fill=white] (8,26-\i*6) rectangle ++(14,-5.5)
      node[midway, font=\scriptsize\ttfamily] {\pc};
  }
  % fetch path
  \node[blk, minimum width=14mm] (icache) at (36,14)
    {\iflangja{命令\\キャッシュ}{I-cache}};
  \node[blk, minimum width=14mm, right=3mm of icache] (dec)
    {\iflangja{デコーダ}{Decoder}};
  \draw[->] (22,14) -- (icache);
  \draw[->] (icache) -- (dec);
  \node[anchor=north, text=ln-muted, align=center] at (36,8.5)
    {\iflangja{選ばれたワープの\\PC で1命令をフェッチ}{one instruction fetched\\at the selected warp's PC}};
  % register file: 4 warps x 4 threads, registers R0..R3
  \node[anchor=south west, font=\scriptsize\sffamily\bfseries] at (66,27)
    {\iflangja{スレッドごとのレジスタ}{a register set per thread}};
  \foreach \w in {0,...,3} {
    \draw[fill=ln-accent!12] (66+\w*13,26) rectangle ++(12.6,-4)
      node[midway] {W\w};
    \foreach \t in {0,...,3} {
      \node[font=\scriptsize\sffamily] at (66+\w*13+\t*3.2+1.6,20.2) {T\t};
    }
  }
  \foreach \r in {0,...,3} {
    \node[anchor=east, inner sep=1pt] at (65.5,16.6-\r*3.4) {R\r};
    \foreach \w in {0,...,3} {
      \foreach \t in {0,...,3} {
        \draw[fill=white] (66+\w*13+\t*3.2,18.3-\r*3.4) rectangle ++(3.2,-3.4);
      }
    }
  }
  \node[anchor=north west, text=ln-muted, align=left] at (66,3)
    {\iflangja{W0 の R1 読み出し = T0--T3 の R1 を同時に読む}{reading R1 for W0 = R1 of T0--T3 at once}};
\end{tikzpicture}
